% =====================================================================
% CAPÍTULO 9: DISCUSIÓN Y ANÁLISIS CRÍTICO
% =====================================================================

\chapter{Discusión y Análisis Crítico}
\label{ch:discusion}

% --- Ajustes editoriales locales del capítulo: sangría, espaciado y control de desbordamientos ---
\setlength{\parindent}{0pt}
\setlength{\parskip}{6pt}
\sloppy
\emergencystretch=2em
\raggedbottom
\section{Introducción a la discusión}
\label{sec:introduccion-discusion}

La discusión interpreta el alcance del proyecto a partir de la literatura revisada y de la problemática institucional. El objetivo no es repetir resultados, sino explicar su sentido técnico y académico, contrastarlos con antecedentes y reconocer límites de la solución propuesta \cite{fieldaware2023,aberdeen2019,gartner2024,garcia2021,rodriguez2022,martinez2020}.

\section{Comparación con soluciones comerciales}
\label{sec:comparacion-comercial}

Las plataformas comerciales ofrecen amplitud funcional y madurez operativa, pero suelen estar diseñadas para configuraciones generales o para empresas con presupuestos más altos. CERMONT S.A.S. presenta una realidad distinta: necesita integrar formatos propios, responder a una operación multiservicio y conservar control sobre sus documentos y procesos. La solución a medida resulta pertinente porque no intenta competir por volumen de funcionalidades, sino por ajuste al contexto.

\section{Aporte del enfoque modular}
\label{sec:aporte-modular}

La arquitectura modular aporta claridad conceptual y técnica. Cada módulo puede entenderse como una respuesta a una necesidad específica del proceso: órdenes, planeación, evidencias, documentos y cierre. Esa organización facilita mantenimiento, pruebas y evolución futura. Desde la literatura de ingeniería de software, esta decisión es consistente con la separación de responsabilidades \cite{bass2021,richards2020,sommerville2016}.

\section{Discusión sobre la trazabilidad documental}
\label{sec:discusion-trazabilidad}

La trazabilidad documental constituye el núcleo del problema y, al mismo tiempo, el núcleo de la solución. Cuando un documento no se enlaza con una orden, una actividad y un responsable, su utilidad operativa disminuye. En cambio, cuando esa relación queda estructurada, el documento se convierte en evidencia útil para el seguimiento y el cierre. Esta observación coincide con la literatura sobre BPM y FSM \cite{dumas2018,fieldaware2023}.

\section{Limitaciones del alcance actual}
\label{sec:limitaciones-discusion}

La solución propuesta tiene límites. Si la versión final no incorpora medición cuantitativa de impacto, no es correcto afirmar reducciones porcentuales. La automatización de formularios desde documentos heredados debe tratarse como línea futura. Además, la arquitectura no debe presentarse como universal sin analizar diferencias en escala, sector o reglas de negocio.

\section{Implicaciones académicas y prácticas}
\label{sec:implicaciones-discusion}

Académicamente, el trabajo muestra cómo una metodología aplicada transforma un problema operativo en un artefacto tecnológico evaluable. Prácticamente, evidencia que una empresa contratista puede organizar mejor su información sin depender exclusivamente de software comercial rígido.

\section{Síntesis crítica}
\label{sec:sintesis-critica}

El aporte del proyecto debe evaluarse desde una posición equilibrada. El aplicativo desarrollado aporta una base técnica para ordenar el flujo documental-operativo de CERMONT S.A.S.; sin embargo, su impacto sobre tiempos, costos, devoluciones o satisfacción de usuarios no debe afirmarse sin medición piloto. Por tanto, la discusión reconoce dos niveles de logro: el logro técnico, respaldado por arquitectura, módulos, capturas y pruebas que deben anexarse; y el logro operativo, pendiente de validarse con usuarios reales y órdenes de trabajo reales.

\section{Discusión sobre la aplicabilidad del modelo a otras organizaciones}
\label{sec:discusion-aplicabilidad}

Aunque el proyecto fue desarrollado para el contexto específico de CERMONT S.A.S., varios elementos de su diseño pueden ser relevantes para organizaciones con perfiles operativos similares. La arquitectura modular, la separación entre el flujo técnico y el flujo administrativo, y el enfoque de formularios dinámicos basados en esquemas declarativos constituyen patrones transferibles a otras empresas contratistas multiservicio, particularmente aquellas que: (a) operan en sectores con alta exigencia documental (hidrocarburos, construcción, mantenimiento industrial); (b) ejecutan servicios en zonas con conectividad variable; (c) manejan una diversidad de formatos operativos que cambian según el tipo de servicio; y (d) requieren trazabilidad de extremo a extremo desde la solicitud del cliente hasta el pago.

Sin embargo, la transferibilidad no es automática. Cada organización tiene formatos propios, reglas de negocio específicas, estructuras organizacionales diferentes y restricciones presupuestales particulares. La arquitectura del sistema fue diseñada para ser adaptable —los formularios se definen mediante esquemas, no mediante código—, pero la adaptación a un nuevo contexto requeriría un diagnóstico documental similar al realizado en este proyecto para identificar los formatos, flujos y puntos de fragmentación propios de la nueva organización. En este sentido, el valor del trabajo no reside en el software como producto terminado y universal, sino en la metodología de diagnóstico, diseño y desarrollo que puede ser replicada en otros contextos.

\section{Discusión sobre el rol de la inteligencia artificial en la evolución del sistema}
\label{sec:discusion-ia}

La incorporación de capacidades de extracción documental basadas en modelos de lenguaje y visión artificial representa una de las líneas de evolución más prometedoras para el sistema. Herramientas como Docling, Unstructured, PaddleOCR y MinerU, revisadas en el estado del arte, demuestran que es técnicamente viable transformar documentos PDF, Word e imágenes en representaciones estructuradas que pueden alimentar un motor de formularios dinámicos. Sin embargo, esta capacidad debe abordarse con cautela académica y técnica.

En primer lugar, la extracción automática de campos desde documentos no estructurados no es un problema resuelto de forma general. La calidad de la extracción depende de la claridad del documento fuente, la consistencia de su formato y la complejidad de su estructura (tablas anidadas, campos multivaluados, secciones condicionales). Los documentos operativos de CERMONT S.A.S., aunque estructurados, presentan variaciones de formato entre diferentes versiones y tipos de servicio que exigirían un proceso de revisión humana antes de publicar cualquier formulario generado automáticamente.

En segundo lugar, desde la perspectiva académica, es esencial no afirmar capacidades que no han sido implementadas ni probadas. La diferencia entre una arquitectura propuesta y una funcionalidad operativa debe mantenerse clara en todo momento. Por esta razón, el módulo de formularios dinámicos desde documentos se presenta en este trabajo como una línea de evolución futura, fundamentada en herramientas existentes y arquitectónicamente compatible con el sistema actual, pero no como un resultado alcanzado.

\section{Discusión sobre la completitud del flujo operativo implementado}
\label{sec:discusion-completitud}

El sistema implementado cubre los módulos principales que abarcan desde la autenticación hasta el cierre administrativo. Sin embargo, la cobertura no es uniforme en profundidad. Los módulos de órdenes de trabajo, evidencias y ejecución alcanzaron un nivel de madurez funcional que incluye operación offline, validación de esquemas y generación de auditoría. Otros módulos, como la aprobación de órdenes de compra (Paso 4) y la firma digital del cliente (Paso 10), tienen la lógica de backend implementada pero carecen de los componentes de interfaz de usuario correspondientes.

Esta distribución asimétrica refleja una decisión consciente de priorización: los módulos que resuelven los cuellos de botella más críticos identificados en el diagnóstico (planeación, ejecución en campo, evidencias, cierre documental) recibieron atención prioritaria, mientras que los módulos de soporte a procesos que actualmente funcionan de forma aceptable para la empresa (gestión de órdenes de compra, firma de actas) se implementaron en el backend para no bloquear el flujo de datos, postergando su interfaz de usuario para iteraciones futuras. Esta estrategia es consistente con el principio de desarrollo iterativo e incremental, que prioriza la entrega de valor temprano sobre la completitud funcional inmediata.

\section{Síntesis de la discusión}
\label{sec:sintesis-final-discusion}

La discusión presentada en este capítulo permite articular una valoración equilibrada del proyecto. Por un lado, el sistema implementado constituye una base técnica documentada, fundamentada en decisiones arquitectónicas, evidencias visuales y pruebas que deben conservarse como anexos. Por otro lado, el trabajo reconoce explícitamente sus limitaciones: la validación cuantitativa de impacto requiere mediciones en entorno real aún no ejecutadas; la cobertura funcional no es uniforme entre módulos; y las capacidades de extracción documental automática, aunque arquitectónicamente previstas, constituyen trabajo futuro. Esta honestidad académica fortalece el trabajo al establecer con precisión qué fue demostrado, qué fue implementado y qué queda como trabajo posterior.

% =====================================================================
\section{Discusión sobre las decisiones arquitectónicas}
\label{sec:discusion-arquitectura}

Las decisiones arquitectónicas documentadas en el marco teórico resultaron pertinentes durante el desarrollo del proyecto. La adopción de un monorepo con npm workspaces ayudó a controlar la duplicación de tipos entre frontend y backend, en coherencia con lo señalado por Richards y Ford \cite{richards2020} sobre la importancia de una fuente única de verdad en arquitecturas distribuidas. La elección de Express 5.2.1 permitió aprovechar la propagación nativa de errores asíncronos y simplificar parte del manejo de excepciones frente a versiones anteriores del framework.

La decisión de utilizar MongoDB con Mongoose (ADR-003) resultó acertada para el contexto de formularios heterogéneos de CERMONT S.A.S., pero introdujo un costo de aprendizaje para desarrolladores acostumbrados a bases de datos relacionales. Este hallazgo coincide con la literatura sobre compensaciones arquitectónicas (\textit{architectural trade-offs}), que señala que ninguna decisión técnica está exenta de contrapartidas \cite{bass2021}.

La implementación de Next.js con App Router (ADR-004) aportó una estructura clara para vistas protegidas y componentes interactivos, aunque la coexistencia de componentes de servidor y cliente exigió una disciplina rigurosa en la delimitación de responsabilidades. La autenticación JWT con Zustand (ADR-005) se consideró adecuada para el escenario offline parcial, pero requiere atención continua a la seguridad de los tokens en el lado del cliente, particularmente en dispositivos compartidos o públicos.

\section{Discusión sobre la cobertura del flujo operativo}
\label{sec:discusion-cobertura}

El sistema implementado cubre módulos principales del flujo operativo de CERMONT S.A.S., desde la autenticación hasta el cierre administrativo. Sin embargo, la cobertura no es uniforme: los módulos de órdenes de trabajo y evidencias cuentan con mayor evidencia funcional, mientras que módulos como la validación automática de certificaciones del personal o la integración directa con portales externos (SAP Ariba) permanecen como líneas de evolución futura.

Esta distribución asimétrica de la madurez funcional es consistente con el enfoque de desarrollo iterativo adoptado, donde los módulos de mayor criticidad para el flujo operativo recibieron prioridad en las iteraciones tempranas. La literatura sobre metodologías ágiles respalda esta estrategia de priorización basada en valor de negocio \cite{beck2001}, aunque también advierte sobre el riesgo de acumular deuda técnica en los módulos postergados.

\section{Discusión sobre la validación del sistema}
\label{sec:discusion-validacion}

La suite de pruebas automatizadas documentada constituye un indicador técnico relevante de la calidad técnica del artefacto desarrollado. Sin embargo, es necesario distinguir entre la validación técnica (que el sistema funcione correctamente) y la validación operativa (que el sistema resuelva el problema de negocio). La primera se sustenta en pruebas automatizadas y debe acompañarse de sus evidencias anexas; la segunda requiere mediciones en un entorno de producción con usuarios reales durante un período suficiente para detectar patrones de uso, fricciones en la interfaz y brechas entre el flujo diseñado y el flujo real de trabajo.

Esta distinción es particularmente relevante cuando el valor académico del proyecto no reside únicamente en el artefacto construido sino en la evidencia de su aplicación efectiva en el contexto organizacional. La literatura sobre Design Science Research enfatiza que la evaluación de un artefacto debe considerar tanto su eficacia técnica (\textit{efficacy}) como su efectividad en el contexto de uso (\textit{effectiveness}) \cite{hevner2004,peffers2007}.

\section{Comparación con trabajos relacionados}
\label{sec:discusion-comparacion}

Los trabajos de grado revisados en el estado del arte abordan problemas análogos de digitalización de procesos técnicos, pero con diferencias significativas en alcance y contexto. García López (2021) desarrolló un CMMS para flota de transporte, logrando resultados en un dominio de mantenimiento de activos propios, pero sin abordar la complejidad documental de una empresa contratista multiservicio. Rodríguez y Torres (2022) implementaron una aplicación móvil para técnicos de campo en telecomunicaciones, priorizando la movilidad sobre la integración administrativa. Martínez Suárez (2020) documentó un sistema web para mantenimiento preventivo industrial, enfocado en la programación de intervenciones más que en el flujo documental completo.

El presente trabajo se diferencia de estos antecedentes en tres aspectos: (a) aborda el ciclo completo desde la solicitud hasta el pago, en lugar de un subconjunto del proceso; (b) integra capacidades offline diseñadas específicamente para operación en campo con conectividad intermitente; y (c) fundamenta sus decisiones técnicas en una revisión sistemática de herramientas comerciales y repositorios abiertos, estableciendo un criterio de comparación verificable.

\section{Lecciones aprendidas del proceso de desarrollo}
\label{sec:lecciones-aprendidas}

El desarrollo del aplicativo dejó varias lecciones relevantes para proyectos similares de Ingeniería Electrónica aplicada a software empresarial. En primer lugar, la adopción temprana de un flujo contract-first (esquema Zod $\rightarrow$ modelo $\rightarrow$ servicio $\rightarrow$ controlador $\rightarrow$ ruta $\rightarrow$ API client $\rightarrow$ hook $\rightarrow$ UI) previno numerosos errores de integración que habrían surgido si el frontend y el backend se hubieran desarrollado de forma independiente. En segundo lugar, la ejecución sistemática de compuertas de calidad después de cada iteración (\texttt{typecheck}, \texttt{lint}, \texttt{test}, \texttt{build}) evitó la acumulación de deuda técnica y facilitó la detección temprana de regresiones. En tercer lugar, la decisión de no incluir datos cuantitativos no verificados en el documento obligó a mantener una disciplina de trazabilidad que, aunque exigente, fortalece la defensa académica del trabajo frente a cualquier revisión académica.

\section{Discusión sobre la expansión del libro a partir de documentación técnica}
\label{sec:discusion-expansion-documentacion-tecnica}

La incorporación de la documentación técnica 00--22 aporta profundidad al trabajo porque permite explicar el desarrollo del aplicativo como proceso y no solo como producto final. Sin embargo, esta integración debe hacerse con criterio académico. No sería adecuado copiar la documentación completa dentro del libro, porque eso transformaría el trabajo de grado en un manual de desarrollo. La estrategia correcta consiste en extraer de cada documento sus decisiones, restricciones, reglas y evidencias, y relacionarlas con los objetivos del proyecto.

Este enfoque también ayuda a responder una posible objeción de jurado: si el libro afirma que se desarrolló un software, debe mostrar cómo se llegó al artefacto. Las secciones ampliadas de metodología, arquitectura, desarrollo, resultados y validación cumplen esa función al convertir la bitácora técnica en argumento académico.

\section{Discusión sobre funcionalidades implementadas y propuestas}
\label{sec:discusion-implementado-propuesto}

Una de las principales fortalezas metodológicas del documento debe ser la separación entre funcionalidad implementada, parcialmente implementada y propuesta. En proyectos de software aplicados a empresas reales, es común que el alcance evolucione y que algunas capacidades queden diseñadas para fases posteriores. Esto no debilita el trabajo si se documenta con transparencia. Por el contrario, demuestra madurez al reconocer límites y al evitar presentar como terminado lo que todavía requiere pruebas o integración.

En el caso de CERMONT, esta distinción es especialmente relevante para módulos como sincronización de evidencias binarias, extracción automática desde documentos, costos reales con datos operativos y validación con usuarios en producción. Es válido incluirlos como arquitectura futura o funcionalidad parcialmente implementada cuando existe base técnica; no es válido reportar impacto cuantitativo sin evidencia.

\section{Discusión sobre arquitectura a medida frente a adopción directa}
\label{sec:discusion-arquitectura-medida-adopcion}

El desarrollo a medida se justifica cuando el problema no se limita a administrar órdenes, sino a integrar un flujo propio de documentos, estados, roles y cierre administrativo. Una herramienta comercial puede cubrir partes del proceso, pero no necesariamente adapta de forma inmediata los documentos internos, la secuencia de SES/facturación, el control de costos reales o las particularidades de evidencias y formatos de la empresa. La decisión de construir una solución propia debe evaluarse por su capacidad de representar fielmente el flujo de CERMONT, no por una preferencia tecnológica.

La discusión también debe reconocer los costos de esta decisión. Un sistema a medida exige mantenimiento, pruebas, documentación, soporte y evolución. Por eso, las recomendaciones deben incluir gobernanza técnica, respaldo de datos, control de versiones, auditoría de seguridad y capacitación de usuarios. Un desarrollo propio sin continuidad puede convertirse en deuda técnica; un desarrollo propio bien documentado puede convertirse en plataforma estratégica.

\section{Lecciones sobre diseño centrado en el proceso}
\label{sec:lecciones-diseno-centrado-proceso}

La principal lección del proyecto es que el software debe diseñarse a partir del proceso real y no desde una lista de pantallas deseadas. En CERMONT, el flujo de valor incluye propuesta, PO, planeación, ejecución, evidencias, informe, acta, SES, factura y pago. Si una pantalla no contribuye a mover una orden dentro de ese flujo o a resolver una falla documentada, su prioridad debe ser revisada. Esta visión evita dashboards decorativos, módulos duplicados y rutas sin propósito.

El diseño centrado en el proceso también facilita explicar la pertinencia del proyecto para Ingeniería Electrónica. La disciplina no se limita a programar interfaces; implica analizar sistemas, modelar estados, controlar transiciones, validar señales de entrada, asegurar integridad, diseñar arquitectura y verificar comportamiento. El aplicativo es un sistema de información aplicado a un proceso técnico-industrial, y por eso puede defenderse dentro del programa académico.

\section{Discusión sobre visualización y comunicación técnica}
\label{sec:discusion-visualizacion-comunicacion}

Las figuras y tablas del libro deben cumplir una función explicativa. Una figura que ocupa una página completa sin texto contextual reduce la fluidez de lectura; una tabla fragmentada artificialmente en varias páginas impide comparar información. Por ello, la corrección editorial debe priorizar diagramas dentro de margen, tablas continuas y explicaciones antes y después de cada elemento visual. No se trata de eliminar figuras, sino de convertirlas en evidencia legible y pertinente.

Los diagramas más importantes para la defensa son: flujo de 14 pasos, mapa de fallas, arquitectura general, modelo de datos, estados de la orden, pipeline de evidencias, generación de informes y validación por capas. Cada uno debe conectarse con un capítulo y con un objetivo. Así, la visualización se convierte en argumento técnico y no en decoración.
